\section{Results and Discussion}
\autoref{tab:phonebench:int-results} presents PR performance, and \autoref{tab:phonebench:ext-results} presents a comprehensive breakdown of downstream evaluations.
In general, \zipactcns performs well in all settings, while \whisper excels in RP.
LALMs generally remain less competitive.

\begin{table*}[tb!]
\centering
\resizebox{0.85\textwidth}{!}{%
\begin{tabular}{ll ccc >{\columncolor{colorint!10}}c ccc >{\columncolor{colorint!10}}c}
\toprule
 & & \multicolumn{4}{c}{\textbf{Variation of Seen Language}} &
 \multicolumn{4}{c}{\textbf{Unseen Languages}} \\
\cmidrule(lr){3-6} \cmidrule(lr){7-10}

\textbf{Model} &  &
{\texttt{PR-tmt} } & {\texttt{PR-arc} } & {\texttt{PR-saa} }  & Avg. & {\texttt{PR-drc} } & {\texttt{PR-vox} } & {\texttt{PR-tsm} } & Avg. \\
\midrule

MultiIPA$^*$          & & 16.3 & 15.5 & 13.8 & 15.2 & 18.3 & 15.2 & 30.5 & 21.3 \\
\lvsixty       & & 13.2 & 10.9 & \wpad9.4 & 11.2 & 17.8 & 15.7 & 24.9 & 19.5 \\
\xlsrf         & & 13.5 & \wpad9.9 & \wpad9.0  & 10.8 & 17.3 & \textbf{13.9} & 31.9 & 21.0 \\
\zipactc       & & \textbf{13.1} & \textbf{\wpad9.7} & \wpad9.0 & \textbf{10.6} & 18.0 & 17.0 & 23.7  & 19.6 \\
\zipactcns     & & \textbf{13.1} & \textbf{\wpad9.7} & \textbf{\wpad8.9} & \textbf{10.6} & \textbf{16.8} & 17.1 & 23.1  & 19.0 \\
\powsm         & & 13.7 & 11.3 & 27.6 & 17.5 & 17.1 & 17.1 & \textbf{22.0} & \textbf{18.7} \\
\powsmctc      & & \textbf{13.1} & 10.3 & 10.0 & 11.1 & 18.1 &  15.3 & 32.2 & 21.9 \\
\midrule
Gemini 2.5 Flash$^{**}$    & & 15.2 & 12.7 & 13.2 & 13.7 & 105.3\wpad & 19.7 & 36.3  & 53.8 \\
Qwen3-Omni-Instruct$^{**}$     & & 15.1 & 11.9 & \wpad9.1 & 12.0 & 150.2\wpad & 49.0 & 117.1\wpad  & 105.4\textcolor{colorint!10}{0} \\

\bottomrule
\end{tabular}
}
\vspace{-5pt}
\caption{PFER of the intrinsic evaluation ($\downarrow$). $^*$English is included during pretraining but not fine-tuning. $^{**}$Some of the ``unseen languages'' may have appeared in the training data. See \autoref{sec:prperf} for details.}
\label{tab:phonebench:int-results}
\vspace{-5pt}
\end{table*}
\begin{table*}[tb!]
\centering
\resizebox{0.95\textwidth}{!}{%
\begin{tabular}{lc lll lll llc >{\columncolor{colorext!10}}c}
\toprule
 & & \multicolumn{3}{c}{\textbf{Pathological Speech}} &
 \multicolumn{3}{c}{\textbf{L2 Speech}} &
 \multicolumn{3}{c}{\textbf{Multilingual Speech}} & 
 \multirow{2}{*}{\shortstack{\textbf{{}} \\ \textbf{Score}}} \\
\cmidrule(lr){3-5} \cmidrule(lr){6-8} \cmidrule(lr){9-11}

\textbf{Model} &  &
{\texttt{DYS-ez} } & {\texttt{DYS-ua} } & {\texttt{CSD-us} } & {\texttt{L1-eda} } & {\texttt{L1-arc} } & {\texttt{L2-so} } & {\texttt{LID-fl} } & {\texttt{GEO-v}} & {\texttt{PI-drc}} &  \\
\midrule

Naive Baseline & & \ms{\wpad0.7}{1.6} & \ms{-0.8}{0.9} & \ms{41.8}{1.0} & \ms{\wpad6.3}{0.4} & \ms{14.3}{0.3} & \ms{\wpad1.5}{1.0} & \ms{\wpad4.3}{0.3} & \ms{\wpad3.3}{0.0} & --- & 8.9 \\
\rowcolor{gray!12}
\addlinespace[0.10cm]
\multicolumn{12}{l}{\textbf{Transcript Probe (TP)}} \\
\addlinespace[0.05cm]
\mipa      & & \ms{48.2}{0.2} & \ms{45.6}{1.4} & \ms{93.6}{1.6} & \textbf{\ms{10.0}{1.0}} & \textbf{\ms{50.5}{0.9}} & \ms{33.3}{1.7} & \ms{89.3}{0.5} & \ms{44.5}{0.4} & 40.9 & \textbf{44.3} \\
\lvsixty   & & \ms{42.4}{1.3} & \ms{50.3}{0.9} & \ms{95.6}{1.4} & \ms{\wpad7.6}{0.5}  & \ms{38.0}{0.3} & \ms{36.1}{1.7} & \ms{91.4}{0.2} & \textbf{\ms{45.7}{0.9}} & 51.3 & 42.0 \\
\xlsrf     & & \ms{49.2}{0.8} & \ms{47.6}{0.8} & \ms{92.3}{2.6} & \wpad\underline{\ms{9.1}{0.6}}  & \underline{\ms{43.1}{0.6}} & \underline{\ms{37.5}{0.8}} & \ms{94.1}{0.2} & \ms{44.5}{1.2} & 56.9 & 43.8 \\
\zipactc   & & \underline{\ms{55.0}{0.6}} & \textbf{\ms{57.0}{0.5}} & \ms{91.7}{2.3} & \ms{\wpad6.6}{0.4}  & \ms{30.5}{0.5} & \ms{36.6}{2.8} & \underline{\ms{95.6}{0.2}} & \ms{44.1}{1.0} & 55.2 & 43.5 \\
\zipactcns & & \textbf{\ms{56.6}{0.8}} & \underline{\ms{51.1}{1.3}} & \textbf{\ms{99.4}{0.5}} & \ms{\wpad6.7}{0.3}  & \ms{30.0}{0.3} & \textbf{\ms{40.8}{0.8}} & \textbf{\ms{95.9}{0.1}} & \underline{\ms{44.7}{1.8}} & \underline{56.6} & \underline{44.2}\\
\powsm     & & \ms{52.7}{1.7} & \ms{46.1}{0.8} & \ms{94.3}{1.3} & \ms{\wpad6.5}{0.8}  & \ms{28.0}{0.3} & \ms{28.4}{2.2} & \ms{95.1}{0.5} & \ms{43.7}{1.4} & 48.7 & 39.6 \\
\powsmctc  & & \ms{53.3}{0.4} & \ms{46.5}{0.6} & \ms{96.9}{0.7} & \ms{\wpad6.4}{0.5}  & \ms{29.8}{0.1} & \ms{26.8}{0.7} & \ms{90.4}{0.4} & \ms{42.9}{0.8} & \textbf{57.7} & 40.2 \\
\midrule
\gemini    & & \ms{27.9}{0.6} & \ms{38.5}{0.4} & \ms{95.0}{1.6} & \ms{\wpad6.4}{0.4}  & \ms{22.3}{0.4} & \ms{20.1}{1.1} & \ms{91.8}{0.3} & \ms{33.2}{1.0} & 39.1 & 31.6 \\
\qweni     & & \ms{52.5}{1.8} & \ms{49.4}{1.0} & \underline{\ms{98.9}{0.8}} & \ms{\wpad6.9}{0.6}  & \ms{30.5}{0.3} & \ms{15.6}{1.4} & \ms{89.3}{0.2} & \ms{34.7}{1.2} & 44.5 & 39.2 \\

\addlinespace[0.10cm]
\midrule
\midrule

\rowcolor{gray!12}
\addlinespace[0.10cm]
\multicolumn{12}{l}{\textbf{Representation Probe (RP)}} \\
\addlinespace[0.05cm]
\mipa       & & \ms{65.5}{4.0} & \ms{77.0}{1.4} & \ms{98.5}{0.8} & \ms{11.7}{1.1} & \ms{53.0}{3.3} & \ms{46.3}{1.9} & \ms{78.2}{1.0} & \ms{24.5}{3.7}  & --- & 56.5 \\
\lvsixty    & & \ms{67.2}{1.8} & \underline{\ms{79.9}{0.9}} & \ms{98.6}{0.6} & \ms{12.0}{0.3} & \ms{60.7}{2.9} & \ms{49.9}{1.0} & \ms{76.6}{1.3} & \underline{\ms{24.6}{2.1}}  & --- & 59.4 \\
\xlsrf      & & \ms{70.8}{2.2} & \textbf{\ms{82.0}{2.2}} & \ms{99.2}{0.7} & \ms{13.0}{1.1} & \ms{47.0}{6.0} & \ms{50.7}{3.1} & \ms{81.0}{2.3} & \ms{21.5}{3.1}  & --- & 58.2 \\
\zipactc    & & \ms{73.2}{2.2} & \ms{74.7}{1.2} & \textbf{\ms{99.5}{0.3}} & \ms{13.9}{0.9} & \ms{73.4}{2.5} & \ms{54.0}{0.8} & \ms{96.1}{0.7} & \ms{23.0}{1.2}  & --- & \underline{62.9} \\
\zipactcns  & & \ms{71.2}{2.2} & \ms{75.1}{0.8} & \ms{98.6}{0.9} & \ms{13.7}{0.6} & \underline{\ms{74.1}{2.8}} & \underline{\ms{54.3}{0.5}} & \textbf{\ms{96.8}{0.3}} & \ms{24.0}{1.5}  & --- & 62.7 \\
\powsm      & & \ms{73.0}{3.0} & \ms{70.8}{1.1} & \textbf{\ms{99.5}{0.3}} & \ms{10.3}{1.2} & \ms{68.0}{1.9} & \ms{53.1}{0.3} & \ms{96.5}{0.1} & \ms{21.5}{2.2}  & --- & 60.4 \\
\powsmctc   & & \underline{\ms{73.6}{1.5}} & \ms{66.7}{1.6} & \ms{97.9}{0.9} & \ms{\wpad8.0}{0.7}  & \ms{53.0}{0.5} & \ms{45.7}{3.0} & \ms{75.4}{1.5} & \ms{14.1}{2.6}  & --- & 55.2 \\
\midrule
\wavlm      & & \ms{69.2}{2.0} & \ms{77.5}{1.4} & \ms{99.0}{0.5} & \underline{\ms{14.4}{1.0}} & \ms{58.3}{2.0} & \ms{50.2}{1.4} & \ms{76.2}{3.2} & \ms{23.5}{4.6}  & --- & 59.4 \\
\whisper    & & \textbf{\ms{74.8}{1.1}} & \ms{79.5}{0.3} & \textbf{\ms{99.5}{0.3}} & \textbf{\ms{24.3}{1.6}} & \textbf{\ms{84.3}{3.0}} & \textbf{\ms{57.2}{0.8}} & \underline{\ms{96.3}{0.5}} & \textbf{\ms{35.0}{2.7}}  & --- & \textbf{68.5} \\

\midrule
\addlinespace[0.10cm]
\rowcolor{gray!12}
\addlinespace[0.10cm]
\multicolumn{12}{l}{\textbf{Zero-shot}} \\
\addlinespace[0.05cm]
\gemini       & & 21.4 & 50.4 & 75.3 & 32.7 & 43.9 & 35.8 & 91.5 & \wpad6.5 & --- & 41.5 \\
\qweni        & & 27.0 & 61.7 & 70.9 & 18.2 & 31.8 & 49.8 & 59.1 & \wpad5.3 & --- & 41.5 \\

\bottomrule
\end{tabular}
}
\vspace{-5pt}
\caption{PR system performance on extrinsic tasks ($\uparrow$). Results are reported as mean $\pm$ standard deviation across 5 random seeds where applicable. Best numbers are \textbf{bolded} and second-best \underline{underlined}. See \autoref{sec:extrinsic} for details.
The formula for aggregrated score is in \autoref{appendix:phonebenchscore}.
}
\vspace{-5pt}
\label{tab:phonebench:ext-results}
\end{table*}

\subsection{Intrinsic Evaluation}
\label{sec:prperf}
We observe a consistent trend for language variation: CTC-based models generally outperform LALMs, followed by AED models. 
For \mipa, English appears during pretraining but not finetuning, highlighting the importance of language coverage in PR data. 
On \texttt{PR-saa}, \powsm performs poorly likely due to decoder search on long speech sequences; meanwhile, a text-based G2P model \cite{zhu2022charsiu-g2p} achieves a PFER of 10.2, beating \gemini despite modeling only canonical pronunciations.

For unseen languages, AED and CTC models show comparable performance, whereas LALMs perform poorly, sometimes producing repeated or degenerate outputs indicative of limited training exposure \cite{neuraltextdegen}. 
Performance also varies across datasets, and language-wise breakdowns reveal heterogeneous behavior. 
\powsm outperforms \powsmctc and exhibits performance comparable to ZIPAs, suggesting that incorporating a degree of language modeling may improve generalization by capturing shared phonological patterns, as further analyzed in \autoref{sec:analysis_noise}.

These trends show that \textbf{variation in seen languages benefits from outputs grounded in known patterns}, whereas \textbf{unseen languages benefit from multilingual training and learned phonological patterns}.


%========================================================================================================================================================

\subsection{Extrinsic Evaluations}
\label{sec:extrinsic}
For transcript probe, ZIPAs and \xlsrf are generally competitive. 
ZIPAs perform well on pathological speech, likely due to their normalized, smaller vocabularies, which approximate broad transcription known to be reliable for speech disorders \cite{shriberg1991reliability}, 
while \xlsrf benefits from diverse pretraining data. 
Multilingual training further improves performance, especially on multilingual tasks.
We discuss \pinv in \autoref{sec:analysis_phoneinv} as an example.
\whisper's strength in representation probe  suggests that large-scale ASR pretraining produces representations that retain phonetic information. 

A trade-off of TP and RP emerges among specialized PR models: for example, Wav2Vec2Phs achieve strong TP results on L2 speech but show limited gains on RP, whereas ZIPAs underperform on TP yet excel on RP. 
Task category also influences their relative performance: Pathological speech benefits more from RP, L2 speech falls in the middle, and multilingual tasks tend to favor TP. 
We hypothesize that transcripts act as a structured bottleneck: pathological speech relies on features such as timbre and prosody, whereas multilingual settings benefit less from acoustic detail.
We investigate the behavior of TP on \geov in \autoref{sec:analysis_geo}.

LALMs show task-dependent performance. Notably, \qweni achieves competitive TP on pathological speech, but they generally perform poorly in zero-shot settings and underperform on languages other than English (\ez). An exception is L2 speech, where the gap is smaller, explored in \autoref{sec:mallm}.

Overall, our results highlight the \textbf{importance of evaluating PR systems with a combination of intrinsic and extrinsic tasks}. 
Intrinsic evaluation alone may not fully capture phonetic capabilities, while extrinsic evaluation reveals that relative performance on TP and RP is task-dependent. 
\textbf{Multilingual pretraining and fine-tuning improve performance} across model families, and \textbf{encoder-CTC based architectures provide more stable PR performance} in new domains. 
In contrast, \textbf{LALMs remain limited} in phone recognition and related tasks.